\section{实验}
\label{sec:experiments}

\subsection{设置与指标}
评价涵盖VideoFeedback~\cite{he2024videoscore}、GenVideo~\cite{chen2024genvideo}
和ComGenVid~\cite{benhayun2026stall}，分别包含11、10和2个生成器单元。
采用STALL的动态等级筛选与生成器内真假平衡配对；
Hotshot-XL和MoonValley保留1秒观察。
共包含15,569个带时长的片段身份、14,573个不同视频文件和23,550条配对记录。
拟合、参考及评价按已知源组隔离；同源的不同观察不作为独立源。
原拟合池每域200片段，其中ComGenVid对应132个源，另两域各200源。

报告AUC及以真实为正类的AP（AP-real），数值越高越好。
每域先对生成器等权平均，本文的总体Average再对三个域等权平均。
对本次匹配预测采用1,000次源组配对bootstrap估计95\%区间；
复用真实视频在不同配对中共享抽样权重。
方法设计在这些开发基准上完成；真实统计拟合不使用生成样本，不等于开发过程从未查看生成结果。

\subsection{主要结果}
\begin{table}[!t]
  \centering
  \caption{锁定 \method 与正式方法链条对比。指标为 pairwise balanced AUC/real-positive AP，数值越高越好。}
  \label{tab:main_results}
  \resizebox{\columnwidth}{!}{%
  \begin{tabular}{lcccc}
    \toprule
    方法 & ComGenVid & VideoFeedback & GenVideo & Macro-3 \\
    \midrule
    Original STALL & 0.8508/0.8560 & 0.8449/0.8596 & 0.8206/0.8129 & 0.8388/0.8428 \\
    Clean single-window & 0.8699/0.8819 & \textbf{0.8631/0.8708} & 0.8382/0.8274 & 0.8570/0.8600 \\
    \method (locked K=3) & \textbf{0.8968/0.9064} & 0.8597/0.8689 & \textbf{0.8657/0.8416} & \textbf{0.8741/0.8723} \\
    \bottomrule
  \end{tabular}
  }
\end{table}

% 预先入队，后续正文按表II的A/B/C和表III顺序引用。
% 由 export_tables.py 生成；CSV 身份及汇总经校验。
\begin{table}[t]
  \centering
  \caption{局部分支的受控贡献。每格为 AUC/AP-real，Average 为三域等权。粗体为整表各列各指标的显示最高值（含并列）；浅底色标出冻结的完整方法。}
  \label{tab:component_ablation}
  {\footnotesize
  \setlength{\tabcolsep}{2.2pt}
  \renewcommand{\arraystretch}{1.15}
  \begin{tabular*}{\columnwidth}{@{\extracolsep{\fill}}lcccc@{}}
    \toprule
    配置 & VF & GV & CGV & Average \\
    \midrule
    \multicolumn{5}{@{}l}{\textbf{A. 分支互补}} \\
    $G$ & 0.823/\textbf{0.844} & 0.873/0.863 & 0.900/0.909 & 0.865/0.872 \\
    $L$（D2） & 0.760/0.774 & 0.851/0.840 & 0.897/0.897 & 0.836/0.837 \\
    \rowcolor{mainred!4}
    $G+L$（Full） & \textbf{0.825}/0.839 & \textbf{0.882}/0.871 & 0.917/0.922 & 0.874/\textbf{0.877} \\
    \addlinespace[3pt]
    \multicolumn{5}{@{}l}{\textbf{B. 先测量后聚合（固定 $G$）}} \\
    pooled D2 & 0.775/0.785 & 0.861/0.851 & \textbf{0.928}/\textbf{0.934} & 0.855/0.857 \\
    D2（同位置数） & \textbf{0.825}/0.839 & \textbf{0.882}/0.872 & 0.917/0.922 & \textbf{0.875}/\textbf{0.877} \\
    \addlinespace[3pt]
    \multicolumn{5}{@{}l}{\textbf{C. 时间与标量替代（固定 $G$）}} \\
    D1 & 0.817/0.828 & 0.878/0.865 & 0.911/0.917 & 0.869/0.870 \\
    TTR & 0.791/0.831 & 0.861/\textbf{0.881} & 0.677/0.736 & 0.777/0.816 \\
    SPLIT 统计 & 0.751/0.803 & 0.814/0.843 & 0.653/0.721 & 0.739/0.789 \\
    二维 Gaussian & 0.778/0.795 & 0.766/0.765 & 0.792/0.802 & 0.779/0.787 \\
    \bottomrule
  \end{tabular*}
  \par\vspace{2pt}
  \begin{minipage}{\columnwidth}\scriptsize
  VF/GV/CGV 为 VideoFeedback/GenVideo/ComGenVid。$G$ 为目标 Global，$L$ 为视频 CDF 后的 Local D2。B、C 固定 $G$ 与 FC 窗口，等权融合，各候选重建 CDF。同位置数指匹配 pooled 的拟合向量数（至多42/片段）。二维 Gaussian 拟合 (TTR, LSMI)；TTR、SPLIT 为同 DINO 控制。
  \end{minipage}}
\end{table}

% 由 export_tables.py 生成；CSV 身份及汇总经校验。
\begin{table}[t]
  \centering
  \caption{更换真实拟合视频后的五次重复。上半为逐域五次平均 AUC/AP-real，Average 为三域等权；下半为总体配对差值均值及95\%区间。逐次重拟合后计算指标，不对预测集成。}
  \label{tab:reference_stability}
  {\footnotesize
  \setlength{\tabcolsep}{2.2pt}
  \renewcommand{\arraystretch}{1.16}
  \begin{tabular*}{\columnwidth}{@{\extracolsep{\fill}}lcccc@{}}
    \toprule
    同池配置 & VF & GV & CGV & Average \\
    \midrule
    $G$ & 0.805/\textbf{0.830} & 0.879/0.871 & 0.915/0.923 & 0.866/0.875 \\
    $G+$D1 & 0.810/0.821 & 0.878/0.867 & 0.925/0.929 & 0.871/0.872 \\
    $G+$D2 & \textbf{0.816}/\textbf{0.830} & \textbf{0.883}/\textbf{0.873} & \textbf{0.929}/\textbf{0.933} & \textbf{0.876}/\textbf{0.879} \\
  \end{tabular*}
  \par\vspace{3pt}
  \begin{tabular*}{\columnwidth}{@{\extracolsep{\fill}}lcc@{}}
    \midrule
    Average 增量 & $\Delta$ AUC [95\% CI] & $\Delta$ AP [95\% CI] \\
    Full$-G$ & \shortstack{$+0.010$\\$[0.005,\,0.015]$} & \shortstack{$+0.004$\\$[-0.001,\,0.009]$} \\
    Full$-$D1融合 & \shortstack{$+0.005$\\$[0.004,\,0.007]$} & \shortstack{$+0.007$\\$[0.005,\,0.008]$} \\
    \bottomrule
  \end{tabular*}
  \par\vspace{2pt}
  \begin{minipage}{\columnwidth}\scriptsize
  Average 样本标准差（AUC/AP）：$G$ 0.005/0.004，$G+$D1 0.004/0.003，$G+$D2 0.004/0.004。
  VF/GV/CGV 同表~\ref{tab:component_ablation}；Full 为 $G+$D2，粗体为上半表各列最高值。每域每池200片段，共2357个新真实视频；新池与旧拟合、评价和参考的已知源组隔离。五池可重叠，区间条件于这五个已选池。
  \end{minipage}}
\end{table}

表~\ref{tab:main_results}报告三个基准上的逐生成器结果。
已有方法的发表值来自STALL的统一比较，本次结果使用目标真实适配及固定预算多窗观察；
视频身份、参考信息及发表Average口径不保证一致，因此两组用于方法背景比较。
本文的归因主要依据下述同协议控制，而不由跨报告差值推断Local贡献。

\method 的Average AUC/AP-real为0.874/0.877。
相对同目标真实片段预算的Global，增量为0.009/0.005，
AUC差值区间为$[0.003,0.016]$，AP区间下界为正但接近零。
两分支共享拟合视频与观察位置，因而该增量不能仅由额外获取目标真实视频解释。
收益具有域差异：VideoFeedback的AP由0.844降至0.839，
而GenVideo和ComGenVid的两项域平均均提高。
